%% ===========================================================================
%%  Bab 1 Pendahuluan
%%  Aturan isi: alinea 4-5 kalimat, satu gagasan per alinea, spasi 2, awal
%%  alinea menjorok otomatis 1,5 cm, rata kiri-kanan. Istilah asing italic:
%%  \asing{open source}.
%% ===========================================================================
\chapter{Pendahuluan}

\section{Latar Belakang}

Kecerdasan buatan berbasis \asing{large language model} telah mengubah cara
peneliti mencari dan menyusun bahan tulisan. Model tersebut mampu merangkum
puluhan halaman dalam hitungan detik dan menawarkan susunan argumen yang belum
terpikirkan oleh penulis. Kemampuan itu membuat model bahasa besar banyak
dipakai sebagai pendamping dalam penyusunan karya ilmiah, mulai dari pencarian
rujukan sampai penyuntingan kalimat. Kemudahan tersebut, bagaimanapun, membawa
persoalan baru yang tidak selalu tampak pada hasil akhir.

Persoalan pertama adalah sitasi yang tidak dapat dipertanggungjawabkan. Model
bahasa besar menyusun jawaban berdasarkan pola yang dipelajarinya, sehingga
rujukan yang dihasilkan dapat tampak meyakinkan tanpa benar-benar merujuk pada
karya yang ada. \citet{ji2023survey} mengelompokkan kekeliruan semacam itu
sebagai halusinasi, yaitu keluaran yang lancar dibaca tetapi tidak sesuai dengan
sumber maupun fakta. \citet{walters2023fabrication} menguji sitasi yang
dihasilkan salah satu model percakapan dan menemukan bahwa sebagian besar di
antaranya tidak akurat, baik karena karyanya tidak ditemukan maupun karena
rincian penulis, tahun, dan terbitannya keliru. Temuan tersebut menunjukkan
bahwa kelancaran bahasa tidak menjamin ketepatan rujukan.

Persoalan kedua menyangkut kendali penulis atas naskahnya sendiri. Alat bantu
menulis yang bekerja langsung pada dokumen dapat mengganti kalimat tanpa
memberi tahu bagian mana yang berubah. Penulis kemudian harus membandingkan
naskah lama dan naskah baru secara manual untuk mengetahui letak perubahannya.
\citet{amershi2019guidelines} menekankan bahwa sistem berbantuan kecerdasan
buatan perlu bersikap terbuka terhadap tindakan pengguna dan menyediakan cara
yang mudah untuk membatalkan perubahan. Tanpa prinsip tersebut, bantuan yang
dimaksudkan mempercepat pekerjaan justru menambah beban pemeriksaan.

Pendekatan yang banyak dipakai untuk menekan halusinasi adalah
\asing{retrieval-augmented generation}. \citet{lewis2020retrieval} merancang
pendekatan ini dengan mengambil dokumen yang relevan terlebih dahulu, lalu
menyusun jawaban berdasarkan dokumen tersebut. \citet{gao2023retrieval}
menjelaskan bahwa mutu jawaban pada pendekatan ini sangat ditentukan oleh
ketepatan tahap pengambilan sumber. Namun, sebagian besar implementasi yang
tersedia masih menempatkan dokumen peneliti pada layanan awan dan belum
membedakan tahap perencanaan dari tahap penulisan naskah. Pemisahan tahap itu
penting agar peneliti dapat menilai rencana tulisan sebelum ada kalimat yang
masuk ke dalam dokumen.

Kebutuhan tersebut muncul secara nyata di Broadband Multimedia Wireless
Laboratory, National Taiwan University of Science and Technology. Laboratorium
ini mengerjakan penelitian di bidang jaringan nirkabel dan multimedia, sehingga
setiap anggota tim mengelola banyak artikel, laporan, dan catatan percobaan
secara bersamaan. Peneliti perlu menelusuri kembali sumber yang pernah dibaca,
memastikan bahwa rujukan yang dipakai benar-benar ada, dan menjaga agar naskah
yang sedang disusun tidak berubah tanpa sepengetahuan mereka. Pekerjaan
penelusuran semacam itu menyita waktu yang seharusnya dapat dipakai untuk
menguji gagasan dan menganalisis hasil.

Berdasarkan keadaan tersebut, selama pelaksanaan magang dirancang dan
diimplementasikan sebuah sistem pencarian data riset berbantuan kecerdasan
buatan yang dijalankan pada komputer peneliti. Sistem ini mencari bagian
dokumen yang relevan dengan pertanyaan peneliti, lalu menyusun jawaban hanya
dari bagian yang ditemukan. Sistem menyediakan tiga cara kerja, yaitu menjawab
pertanyaan, menyusun kerangka, dan menyusun draf bagian naskah. Setiap usulan
perubahan harus disetujui peneliti sebelum dimasukkan ke dalam naskah, dan
setiap sitasi diperiksa terhadap pustaka yang tersimpan pada proyek.

Sistem yang dibangun juga mencatat asal setiap potongan sumber sehingga jawaban
dapat ditelusuri kembali. Pencatatan ini penting karena peneliti tidak cukup
mempercayai jawaban yang terbaca lancar; mereka perlu mengetahui dari dokumen
mana sebuah pernyataan berasal. Naskah yang telah disetujui diekspor dalam
format LaTeX beserta daftar pustakanya, lalu disinkronkan dengan repositori
proyek pada GitHub. Dengan demikian, riwayat perubahan naskah tetap terekam dan
dapat diperiksa oleh pembimbing maupun rekan tim.

Laporan magang ini merekam perancangan, implementasi, dan evaluasi sistem
tersebut sebagai bagian dari pelaksanaan magang pada periode 12 Januari sampai
28 September 2026. Laporan disusun mengikuti Pedoman Penulisan Laporan Magang
STMIK Kuwera agar dapat dipertanggungjawabkan secara akademik. Susunan bab
disesuaikan dengan urutan yang ditetapkan pedoman, mulai dari dasar pemikiran
sampai kesimpulan. Selain menjadi dokumen pertanggungjawaban, laporan ini
diharapkan menjadi rujukan bagi mahasiswa yang akan mengerjakan pekerjaan
serupa pada periode berikutnya.

\section{Tujuan Magang}

Magang dilaksanakan untuk menghubungkan pengetahuan yang diperoleh di bangku
kuliah dengan pekerjaan nyata di lingkungan laboratorium riset. Penulis
ditempatkan pada pekerjaan pengembangan perangkat lunak yang menopang kegiatan
penelitian, sehingga dapat mengamati langsung bagaimana kebutuhan peneliti
diterjemahkan menjadi fitur sistem. Penempatan tersebut juga memberi kesempatan
kepada penulis untuk menguji keputusan rancangan pada pengguna yang sebenarnya.
Tujuan umum tersebut dijabarkan menjadi beberapa tujuan khusus sebagai berikut.

\begin{enumerate}
  \item Merancang sistem pencarian data riset berbantuan kecerdasan buatan yang
        dijalankan secara lokal pada komputer peneliti.
  \item Mengimplementasikan pencarian dokumen sumber sehingga jawaban yang
        dihasilkan dapat ditelusuri kembali kepada bagian asalnya.
  \item Mengimplementasikan verifikasi sitasi terhadap pustaka lokal yang
        tersimpan pada setiap proyek penelitian.
  \item Mengimplementasikan alur revisi yang menempatkan persetujuan peneliti
        sebagai syarat sebelum naskah berubah.
  \item Menguji ketepatan sumber, ketepatan penentuan bagian naskah, serta
        keberhasilan ekspor dan sinkronisasi naskah.
\end{enumerate}

Melalui kelima tujuan tersebut, magang diarahkan pada satu hasil yang utuh,
yaitu perangkat lunak yang dapat dipakai dalam pekerjaan penulisan sehari-hari.
Setiap tujuan memiliki ukuran keberhasilan yang diperiksa pada tahap pengujian.
Penulis juga mencatat kendala yang muncul selama pengerjaan agar dapat dijadikan
bahan perbaikan pada pengembangan berikutnya.

\section{Manfaat Magang}

Pelaksanaan magang memberikan manfaat bagi tiga pihak yang terlibat, yaitu
peserta magang, program studi, dan instansi tujuan magang. Manfaat tersebut
tidak berdiri sendiri, melainkan saling berkaitan sebagai hasil dari pekerjaan
yang sama. Peserta magang memperoleh pengalaman kerja, program studi memperoleh
contoh penerapan ilmu, dan instansi memperoleh perangkat yang membantu
pekerjaannya. Uraian berikut menjelaskan bentuk manfaat pada masing-masing
pihak.

\subsection{Manfaat bagi Peserta Magang}

Manfaat utama bagi peserta magang adalah pengalaman mengerjakan perangkat
lunak dari tahap perancangan sampai pengujian. Penulis belajar menerjemahkan
kebutuhan pengguna yang masih umum menjadi rancangan fitur yang dapat
dikerjakan. Penulis juga berlatih menulis kode, membaca kesalahan, dan
memperbaiki sistem yang sudah berjalan tanpa mengganggu data pengguna.
Pengalaman tersebut tidak diperoleh hanya dari membaca buku, karena setiap
keputusan harus diuji pada pemakaian yang sebenarnya.

Selain keterampilan teknis, magang melatih cara berkomunikasi dengan peneliti
yang menjadi pengguna sistem. Penulis belajar menanyakan kebutuhan dengan
pertanyaan yang tepat dan menjelaskan keputusan teknis tanpa memakai istilah
yang membingungkan. Pengalaman ini menumbuhkan kebiasaan mendokumentasikan
pekerjaan agar dapat dilanjutkan oleh orang lain. Penulis juga memperoleh
gambaran mengenai budaya kerja penelitian di laboratorium universitas luar
negeri.

\subsection{Manfaat bagi Program Studi}

Bagi Program Studi Sistem Informasi, magang ini menghasilkan contoh penerapan
ilmu sistem informasi pada pekerjaan penelitian. Kasus yang dikerjakan
memperlihatkan bahwa analisis kebutuhan, perancangan basis data, dan
pengembangan antarmuka tidak hanya dipakai pada sistem bisnis, tetapi juga pada
perangkat pendukung riset. Contoh tersebut dapat dipakai sebagai bahan
pembelajaran pada mata kuliah yang relevan. Dengan begitu, materi kuliah
memperoleh gambaran penerapan yang dekat dengan pekerjaan lulusan.

Magang ini juga memperluas jejaring kerja sama program studi dengan institusi
pendidikan di luar negeri. Hubungan yang terbentuk selama periode magang dapat
dimanfaatkan untuk kegiatan pertukaran mahasiswa maupun kerja sama penelitian
pada masa berikutnya. Pengalaman penulis selama mengikuti program tersebut
dapat dibagikan kepada mahasiswa lain sebagai gambaran dunia kerja di bidang
tersebut. Program studi memperoleh masukan mengenai kemampuan yang perlu
diperkuat pada mahasiswa yang berminat bekerja di bidang riset.

\subsection{Manfaat bagi Instansi Tujuan Magang}

Bagi Broadband Multimedia Wireless Laboratory, magang ini menghasilkan
perangkat lunak yang membantu anggota laboratorium mengelola bahan penelitian.
Sistem yang dibangun menyimpan dokumen sumber pada komputer peneliti dan
memungkinkan penelusuran kembali bagian yang relevan dengan pertanyaan. Dengan
demikian, waktu yang biasa dipakai untuk mencari ulang dokumen dapat dialihkan
ke pekerjaan analisis. Perangkat tersebut juga dapat dipakai oleh anggota baru
yang sedang mempelajari topik penelitian laboratorium.

Laboratorium juga memperoleh alur kerja yang mencatat persetujuan peneliti
sebelum naskah berubah. Pencatatan tersebut memudahkan peninjauan bersama antara
anggota tim dan pembimbing, karena setiap perubahan dapat ditelusuri. Selain
itu, sistem menyimpan naskah beserta daftar pustakanya pada repositori proyek,
sehingga riwayat penyusunan dokumen tetap tersedia setelah mahasiswa magang
menyelesaikan periodenya. Hasil kerja ini menjadi titik awal bagi pengembangan
perangkat serupa pada proyek penelitian berikutnya.

